\documentclass[10pt,a4paper]{amsart}
\usepackage[margin=19mm]{geometry}
\usepackage{amsmath,amssymb,mathtools}
\usepackage[scr=rsfs,cal=euler]{mathalfa}
\usepackage{xfrac}
\usepackage[unicode,hidelinks,bookmarks=false]{hyperref}
\newcommand{\ie}{i.e.\ }
\newcommand{\cf}{cf.\ }
\newcommand{\resp}{respectively\ }
\newcommand{\Subsection}{\subsection}
\providecommand{\nobreakdash}{\nobreak\hbox{-}}
\setlength{\emergencystretch}{2em}
\multlinegap0pt
\usepackage[shortlabels]{enumitem}
\usepackage{amssymb}
\usepackage{tikz}
\newtheorem{theorem}{Theorem}
\newtheorem{proposition}{Proposition}
\newcommand{\R}{\mathbb{R}}
\title{On semilinear Grushin--Schrödinger equation in $\R^N$}
\author{J\^onison Carvalho, Arl\'ucio Viana}
\date{Source: arXiv:2603.06417v1 (CC BY 4.0)}
\begin{document}
\maketitle

\noindent\emph{Source and licence.} On semilinear Grushin--Schrödinger equation in $\R^N$. Version arXiv:2603.06417v1 (CC BY 4.0), distributed by arXiv under the Creative Commons Attribution 4.0 International licence, http://creativecommons.org/licenses/by/4.0/, as recorded in the arXiv licence metadata for this exact version and shown on the abstract page https://arxiv.org/abs/2603.06417v1. Full licence notice: \texttt{licenses/arxiv-2603.06417-CC-BY-4.0.txt}.\par\smallskip

\noindent\textit{Editorial scope.} Counted unit: the article's proposition proposition (source0.tex lines 626--628). The article's complete original proof of it is delivered with it (source0.tex lines 630--785, one \texttt{proof} environment of 156 lines). No mathematics is written, repaired or re-derived: the statement and the proof are the article's own lines, in the article's own notation, and no retained line is substituted or re-flowed.  Retained uncounted as the source-local context the proof needs: lines 136--146; lines 137--139; lines 148--153; lines 203--236; lines 249--252; lines 261--267; lines 287--295; lines 487--489; lines 604--607.  Nothing was omitted from any retained span, so the package carries no omission record.  No retained line contains a citation, so the wrapper carries no bibliography and no bibliography entry.  No retained line contains a figure environment, an image-inclusion command or drawing code, so no image is delivered and the package declares no asset dependency.  Editorial formatting only, in addition: an AMS \texttt{amsart} preamble that restates the article's own declarations and macros used by the retained lines, namely the environments \texttt{theorem}, \texttt{proposition} on separate counters, together with the standard packages those lines need.  The retained environments print in this document as Theorem 1, Proposition 1 in order of appearance, whereas the article prints the same items with the numbering of its own full document.  The complete native source of the article as distributed in the arXiv e-print for this exact version is bundled unchanged as \texttt{sources/195869/source0.tex}.
The main goal of the present work is to establish the compact embedding $$E_V^\gamma \hookrightarrow L^p_Q(\mathbb{R}^N) ,$$ under suitable conditions on $V$ and $Q$. As consequence, we prove the existence of nonnegative weak solutions to the equation
\begin{equation}\tag{$\mathcal{P}$}\label{P}
	-\Delta_\gamma u + V(z)u = Q(z)f(u), \quad z\in \mathbb{R}^N.
\end{equation}
In this work, we assume instead the following assumptions on the potentials:
\begin{enumerate}[label=($V$),ref=$(V)$]
	\item \label{V} $V: \mathbb{R}^N\to \mathbb{R}$ is a measurable function and there exist constants  $C_V>0$ and $\alpha\in\mathbb{R}$ such that
	$$
	\frac{C_V}{\left(1+|z|\right)^\alpha}\leq V(z), \quad\mbox{for a.e. } z\in\mathbb{R}^N;
	$$
\end{enumerate}

\begin{equation}\tag{$\mathcal{P}$}
	-\Delta_\gamma u + V(z)u = Q(z)f(u), \quad z\in \mathbb{R}^N.
\end{equation}

\begin{enumerate}[label=($Q$),ref=$(Q)$]
	\item \label{Q} $Q: \mathbb{R}^N\to \mathbb{R}$ is a measurable function and there there exist constants   $C_Q>0$ and $\beta\in\mathbb{R}$ such that
	$$
	0<Q(z)\leq \frac{C_Q}{\left(1+|z|\right)^\beta}, \quad\mbox{for a.e. } z\in\mathbb{R}^N.
	$$
\end{enumerate}

\begin{theorem}[Weighted Sobolev Embedding]
	\label{imersao} Assume that \ref{V} and \ref{Q} hold. Then the embedding $E_V^\gamma \hookrightarrow L^p_Q(\mathbb{R}^N)$ is continuous  if one of the following conditions holds: 	
	\begin{itemize}
		\item [(i)] $\alpha \leq  N < \beta$ and $2\leq p \leq 2_\gamma^*$;
		
		\item[(ii)] $N
        \leq\alpha < \beta$ and $2\leq p \leq \overline{p}$,
		where
		$$
		\overline{p}:=\min\left\lbrace \frac{2(\beta-N)}{(\alpha-N)},2_\gamma^* \right\rbrace;
		$$
		
		\item [(iii)] $\alpha\leq N$, $ \underline{\beta} \leq \beta <N$, and $\underline{p}\leq p\leq2_\gamma^*$,
		where
		$$
		\underline{\beta}:=\max\left\lbrace N-\frac{2_\gamma^*(N-2)}{2},N-\frac{2_\gamma^*(N-\alpha)}{2}\right\rbrace 
		$$
		and
		$$
		\underline{p}:=\max\left\lbrace 2, \frac{2(N-\beta)}{(N-2)}, \frac{2(N-\beta)}{(N-\alpha)} \right\rbrace.
		$$
		
		
	\end{itemize}
	Furthermore, the embedding is compact whenever:
	\begin{itemize}
		\item [(i')] $\alpha \leq  N < \beta$ and $2\leq p < 2_\gamma^*$;
		
		\item[(ii')] $N<\alpha\leq\beta$ and $2\leq p < \overline{p}$.
		
		
		%\item [(iii')] $\alpha<N$, $ \underline{\beta} \leq \beta <N$, and $\underline{p}< p<2_\gamma^*$.
	\end{itemize}
\end{theorem}

\begin{enumerate}[label=($f_1$),ref=$(f_1)$]
	\item \label{(f_1)} 
    \(\displaystyle\lim_{s\to 0}\dfrac{f(s)}{s}=0\);
\end{enumerate}

\begin{enumerate}[label=($f_3$),ref=$(f_3)$]
	\item \label{(f_3)} there exists $\theta>2$ such that
	$$
	0<\theta F(s) \leq f(s) s, \quad \forall \, s \in \mathbb{R} \backslash\{0\},
	$$	
	where $F(t):=\int_0^t f(s) \,\mathrm{d} s$.
\end{enumerate}

\begin{theorem}\label{existencia_de_solucao} Let the assumptions of the items $(i')$ and $(ii')$ in Theorem \ref{imersao} hold, and assume \ref{V}, \ref{Q}, and \ref{(f_1)}-\ref{(f_3)}. Then problem \eqref{P} has at
	least one nonnegative nontrivial weak solution $u\in E_V^\gamma$. Furthermore,
	\begin{itemize}
		\item [$(i)$] if $\dfrac{1}{Q} \in L^\infty_{\mathrm{loc}}(\R^N)$, then $u\in L^\infty_{\mathrm{loc}}(\mathbb{R}^N)$;

        
		\item[$(ii)$] if there exist $V_0, V_1>0$ such that $V_0\leq V(z)\leq V_1$, for all $z\in\R^N$, then $u \in W_\gamma^{2,\frac{2^*_\gamma}{q-1}}(\R^N)$.
	\end{itemize}
\end{theorem}

\begin{equation}\label{fcond}
	|f(s)| \leq \varepsilon|s| + C|s|^{q-1}	\quad \mbox{and} \quad |F(s)| \leq \varepsilon s^2 + C|s|^{q}, \quad \forall \, s \in \mathbb{R}.
\end{equation}

    \begin{proposition}\label{existencia-solucao}
	Let the assumptions of Theorem \ref{existencia_de_solucao} hold. Then problem \eqref{P} has at
	least one nonnegative nontrivial weak solution.
\end{proposition}

    \begin{proposition}
	Let the assumptions of Theorem \ref{existencia_de_solucao} hold. If $\dfrac{1}{Q}\in L^\infty_{\mathrm{loc}}(\R^N)$, then the solution obtained in Proposition~\ref{existencia-solucao} belongs to $L^\infty_{\mathrm{loc}}(\mathbb{R}^N)$.
\end{proposition}

\begin{proof}
	For each $n \in \mathbb{N}$ and $\xi>1$, we consider the sets $A_n:=\left\lbrace z \in \mathbb{R}^N : u^{\xi-1}(z) \leq n\right\rbrace $ and $D_n:= \mathbb{R}^N \setminus A_n$. 
Let
	$$
	w_n:=\left\{
	\begin{aligned}
		& u^{2\xi-1} ,\ &\mbox{in} \ A_n, \\
		&n^2u, \ &\mbox{in} \ D_n.
	\end{aligned}
	\right.
	$$
We claim that \(w_n \in E_V^\gamma\). In fact, see that
	$$
	\nabla_{\gamma}w_n=\left\{
	\begin{aligned}
		& (2\xi-1)u^{2(\xi-1)}\nabla_{\gamma}u ,\ &\mbox{in} \ A_n, \\
		&n^2\nabla_{\gamma}u, \ &\mbox{in} \ D_n.
	\end{aligned}
	\right.
	$$
    This implies
	$$
	\int_{\R^N}|\nabla_{\gamma} w_n|^2\,\mathrm{d}z \leq (2\xi-1)^2n^4\int_{A_n}|\nabla_\gamma u|^2\,\mathrm{d}z + n^4\int_{D_n}|\nabla_\gamma u|^2\,\mathrm{d}z<\infty
	$$
	and
\[
\begin{aligned}
    \int_{\R^N}V(z)w_n^2\,\mathrm{d}z &= \int_{A_n}V(z)u^{4(\xi-1)}u^2\,\mathrm{d}z + n^4\int_{D_n}V(z)u^2\,\mathrm{d}z 
    \\
    &\leq n^4\int_{\R^N}V(z)u^2\,\mathrm{d}z<\infty.
\end{aligned}
\]
Consequently, $w_n \in E_V^\gamma$ and the claim is proved.

On other hand, since $u$ is solution of \eqref{P}, then
\begin{equation}\label{solucao}
		\int_{\mathbb{R}^N}\nabla_\gamma u \cdot\nabla_\gamma w_n\,\mathrm{d}z + \int_{\mathbb{R}^N}V(z)uw_n\,\mathrm{d}z 		= \int_{\R^N}Q(z)f(u)w_n\,\mathrm{d}z 	.
	\end{equation}
    Now, consider
	$$
	v_n:=\left\{
	\begin{aligned}
		& u^{\xi} ,\ &\mbox{in} \ A_n, \\
		&nu, \ &\mbox{in} \ D_n.
	\end{aligned}
	\right.
	$$
	Hence, 
	$$
	\nabla_\gamma v_n:=\left\{
	\begin{aligned}
		& \xi u^{\xi-1}\nabla_\gamma u ,\ &\mbox{in} \ A_n, \\
		&n\nabla_\gamma u, \ &\mbox{in} \ D_n.
	\end{aligned}
	\right.
	$$
	Similarly, $v_n \in E_V^\gamma$. Notice that
	\begin{equation}\label{derivada-fraca-wn}
		\int_{\mathbb{R}^N}\nabla_\gamma u \cdot \nabla_\gamma w_n\,\mathrm{d}z = (2\xi-1)\int_{A_n}u^{2(\xi-1)}|\nabla_{\gamma} u|^2\,\mathrm{d}z + n^2\int_{D_n}|\nabla_\gamma u|^2\,\mathrm{d}z
	\end{equation}
	and
    $$
	\int_{\mathbb{R}^N}|\nabla_{\gamma} v_n|^2\,\mathrm{d}z = \xi^2\int_{A_n}u^{2(\xi-1)}|\nabla_{\gamma} u|^2\,\mathrm{d}z + n^2\int_{D_n}|\nabla_\gamma u|^2\,\mathrm{d}z.
	$$
So,
	$$
	\int_{\mathbb{R}^N}|\nabla_{\gamma} v_n|^2\,\mathrm{d}z - \int_{\mathbb{R}^N}\nabla_\gamma u \cdot \nabla_\gamma w_n\,\mathrm{d}z = (\xi^2-2\xi+1)\int_{A_n}u^{2(\xi-1)}|\nabla_{\gamma} u|^2\,\mathrm{d}z.
	$$
From this and \eqref{derivada-fraca-wn}, it holds that
	$$
	\begin{aligned}
		\int_{\mathbb{R}^N}|\nabla_{\gamma} v_n|^2\,\mathrm{d}z &\leq \frac{(\xi^2-2\xi+1)}{(2\xi-1)}\int_{\mathbb{R}^N}\nabla_\gamma u \cdot \nabla_\gamma  w_n\,\mathrm{d}z + \int_{\mathbb{R}^N}\nabla_\gamma u \cdot \nabla_\gamma w_n\,\mathrm{d}z
		\\
		& < \xi^2\int_{\mathbb{R}^N}\nabla_\gamma u \cdot \nabla_\gamma w_n\,\mathrm{d}z.
	\end{aligned}
	$$
Applying \eqref{solucao},
    $$
	\int_{\mathbb{R}^N}|\nabla_{\gamma} v_n|^2\,\mathrm{d}z \leq \xi^2\left(	 \int_{\R^N}Q(z)f(u)w_n\,\mathrm{d}z - \int_{\mathbb{R}^N}V(z)uw_n\,\mathrm{d}z 	\right).
	$$	
Combining the above inequality with the fact that $\xi>1$ and the equality
	$$
	\int_{\mathbb{R}^N}V(z)v_n^2\,\mathrm{d}z = \int_{\mathbb{R}^N}V(z)uw_n\,\mathrm{d}z,
	$$	
	we infer that
\begin{equation}\label{est-f(u)}
		\|v_n\|_{E_V^\gamma}^2 <\xi^2	 \int_{\R^N}Q(z)f(u)w_n\,\mathrm{d}z.
	\end{equation}
    From \eqref{fcond} and the estimate $w_n \leq u^{2\xi-1}$ in $\mathbb{R}^N$, one deduces
$$
	\int_{\R^N}Q(z)f(u)w_n\,\mathrm{d}z \leq \varepsilon\int_{\mathbb{R}^N}Q(z)u^{2\xi}\,\mathrm{d}z + C\int_{\mathbb{R}^N}Q(z)u^{2\xi+q-2}\,\mathrm{d}z.
	$$ 
    By H\"older's inequality,
	$$
	\begin{aligned}
		\int_{\R^N}Q(z)f(u)w_n\,\mathrm{d}z &\leq \varepsilon\left(\int_{\mathbb{R}^N}Q(z)\,\mathrm{d}z\right)^{\frac{1}{r_1}}\left(\int_{\mathbb{R}^N}Q(z)u^{2r_2\xi}\,\mathrm{d}z\right)^{\frac{1}{r_2}}
		\\
		&\quad + C\left(\int_{\mathbb{R}^N}Q(z)u^{2q_1\xi}\,\mathrm{d}z\right)^{\frac{1}{q_1}}\left(\int_{\mathbb{R}^N}Q(z)u^{q_2(q-2)}\,\mathrm{d}z\right)^{\frac{1}{q_2}},
	\end{aligned}
	$$	
	where $\frac{1}{r_1} + \frac{1}{r_2} =  \frac{1}{q_1} + \frac{1}{q_2}=1$, with 
$$
	\frac{2}{q-2}\leq q_2-1< q_2< \frac{2_\gamma^*}{q-2}.
	$$
	This is possible because $\frac{2_\gamma^*}{q-2} - \frac{2}{q-2}>1$, for any $q\in(2,2_\gamma^*)$. Since $\beta>N$, $Q\in L^1(\mathbb{R}^N)$. Moreover, $E_V^\gamma \hookrightarrow L_Q^{q_2(q-2)}(\mathbb{R}^N)$ (see Theorem~\ref{imersao}). Hence, if $q_1=r_2$, then
	$$
	\int_{\R^N}Q(z)f(u)w_n\,\mathrm{d}z \leq C_1\|u\|_{L^{2q_1\xi}_Q(\mathbb{R}^N)}^{2\xi},
	$$
	for some $C_1=C_1(\|u\|_{E^\gamma_V})>0$. This and \eqref{est-f(u)} infer that 
	$$
	\|v_n\|_{E_V^\gamma}^2 <\xi^2C_1\|u\|_{L^{2q_1\xi}_Q(\mathbb{R}^N)}^{2\xi}.
	$$
By Theorem~\ref{imersao} and the last inequality, we have
	$$
\left(\int_{A_n}Q(z)u^{2_\gamma^*\xi}\,\mathrm{d}z\right)^{\frac{2}{2_\gamma^*}}\leq\|v_n\|_{L^{2_\gamma^*}_Q(\mathbb{R}^N)}^2 \leq \xi^2C_2\|u\|_{L^{2q_1\xi}_Q(\mathbb{R}^N)}^{2\xi}.
	$$
where we used that $v_n=u^\xi$ in $A_n$. Taking the limit as $n\to \infty$ and using Fatou’s lemma, we reach
	$$
	\|u\|_{L^{2_\gamma^*\xi}_Q(\mathbb{R}^N)}^{2\xi} \leq \xi^2C_2\|u\|_{L^{2q_1\xi}_Q(\mathbb{R}^N)} 
	$$
    or
	\begin{equation}\label{inicial-ite-moser}
		\|u\|_{L^{2_\gamma^*\xi}_Q(\mathbb{R}^N)} \leq \xi^{\frac{1}{\xi}} C_2^{\frac{1}{2\xi}}\|u\|_{L^{2q_1\xi}_Q(\mathbb{R}^N)}.
	\end{equation}
Now, we start Moser's iteration, taking $\xi=\sigma= \frac{2_\gamma^*}{2q_1}>1$. Indeed,
	$$
	q_1=\frac{q_2}{q_2-1}=q_2\cdot\frac{1}{q_2-1}<\frac{2_\gamma^*}{q-2}\cdot \frac{q-2}{2} = \frac{2_\gamma^*}{2}.
	$$
Thus,
	$$
\|u\|_{L^{2_\gamma^*\sigma}_Q(\mathbb{R}^N)} \leq \sigma^{\frac{1}{\sigma}} C_2^{\frac{1}{2\sigma}}\|u\|_{L^{2_\gamma^*}_Q(\mathbb{R}^N)}.
	$$
    Taking $\xi=\sigma^2$ in \eqref{inicial-ite-moser} and noting that $2q_1\sigma^2=(2q_1\sigma)\sigma = 2_\gamma^*\sigma$, the last estimate, yields
	$$
	\|u\|_{L^{2_\gamma^*\sigma^2}_Q(\mathbb{R}^N)} \leq (\sigma^2)^{\frac{1}{\sigma^2}} C_2^{\frac{1}{2\sigma^2}}\|u\|_{L^{2_\gamma^*\sigma}_Q(\mathbb{R}^N)} \leq \sigma^{\frac{1}{\sigma}+\frac{2}{\sigma^2}}C_2^{\frac{1}{2\sigma} + \frac{1}{2\sigma^2}}\|u\|_{L^{2_\gamma^*}_Q(\mathbb{R}^N)}.
	$$
By continuing the iteration, we can conclude that
	$$
	\|u\|_{L^{2_\gamma^*\sigma^j}_Q(\mathbb{R}^N)} \leq \sigma^{\sum_{k=1}^\infty \frac{k}{\sigma^k}} C_2^{\sum_{k=1}^\infty \frac{1}{2\sigma^k}}\|u\|_{L_Q^{2_\gamma^*}(\mathbb{R}^N)}.
	$$
Since these series converge, we obtain
	$$
	\|u\|_{L^{2_\gamma^*\sigma^j}_Q(\mathbb{R}^N)} \leq C_3\|u\|_{L_Q^{2_\gamma^*}(\mathbb{R}^N)}.
	$$


    Now, consider $R>0$ arbitrary. From the hypothesis $\dfrac{1}{Q} \in L^\infty_{\mathrm{loc}}(\R^N)$, there exists $Q_R>0$ such that $Q(z) \geq Q_R$, for any $z \in B_R$. Thus,
	$$
	\int_{B_R}|u|^{2_\gamma^*\sigma^j} \mathrm{d}z \leq \frac{1}{Q_R}\int_{B_R}Q(z)|u|^{2_\gamma^*\sigma^j} \mathrm{d}z \leq \frac{1}{Q_R}\|u\|_{L^{2_\gamma^*\sigma^j}_Q(\mathbb{R}^N)}^{2_\gamma^*\sigma^j} \leq \frac{C_3^{2_\gamma^*\sigma^j}}{Q_R}\|u\|_{L^{2_\gamma^*}_Q(\mathbb{R}^N)}^{2_\gamma^*\sigma^j}.
	$$
Consequently,
	$$
	\|u\|_{L^{2_\gamma^*\sigma^j}(B_R)} \leq \frac{C_3}{Q_R^{\frac{1}{2_\gamma^*\sigma^j}}}\|u\|_{L_Q^{2_\gamma^*}(\mathbb{R}^N)}.
	$$
Taking $j \to \infty$, it holds that $u \in L^\infty(B_R)$, finalizing the proof.
\end{proof}

\end{document}
